\documentclass{article}
\usepackage[T1]{fontenc}
\usepackage{amsmath,amssymb}
\begin{document}
\section*{Dependent Fubini: what the binary kernel establishes}
\paragraph{Statement.}
Fix two independent occurrences and a target pair of parent and continuation
addresses at each. The source of the comparison is the explicitly ordered
sum over $(x,y,u,v)$, with parent indices $x,y$ and continuation indices
$u\in B_x$, $v\in D_y$. The coefficient at a leaf is
$(V_x\otimes B_{xu})\otimes(W_y\otimes D_{yv})$.
The outer source lists need not be nonempty, the inner lists may be ragged,
and equal coefficient objects at different positions do not identify those
positions. Given arbitrary arrows
\[
 \beta_x:\bigoplus_u B_{xu}\longrightarrow b,\quad
 \theta_y:\bigoplus_v D_{yv}\longrightarrow d,\quad
 \alpha:\bigoplus_x V_x\longrightarrow a,\quad
 \gamma:\bigoplus_y W_y\longrightarrow c ,
\]
\texttt{ragged\_binary} equates two recursively constructed arrows to
$(a\otimes c)\otimes(b\otimes d)$. The separate schedule shuffles after
the parent merges. The joint schedule shuffles the four factors at every
leaf, then merges the continuation resources, then the parents.
\paragraph{Structural input.}
Addition and multiplication act bifunctorially on arrows. Gathering uses
binary distributors and nullary absorbers. The tensor middle-four is
natural. Its compatibility with these gatherings is tested at one binary
sum and one zero in each of the four slots. These eight local diagrams
contain no finite grids, target fibres, native arrows or desired final
equality. Folding each local diagram gives the two rectangle comparisons.
The general binary result is conditional on those diagrams: their status
cannot be inferred from the names of two bifunctors. The reverse-function
sum/product rig discharges every one. The mathlib adapter recovers the genuine
distributor/absorber inverses, additive structure, braided middle, and relevant
naturality from the repository's grid API; that API by itself does not
supply all eight diagrams from the standard rig presentation.
\paragraph{Native arrows and augmentation.}
Finite addresses are retained in the preimage of the actual dependent map
$(x,u)\mapsto(fx,g_xu)$. In the abstract presentation a coherent additive
direct image is explicit input. The category and slice over the one-address
unit family then have proved identities, associativity and terminal object.
Stable partition of finite lists is constructed from additive units,
associators and swaps; its naturality is proved separately, not its entire
comparison with the chosen direct image. The two-stage theorem proves
interchange across arbitrarily many intermediate parent blocks for a common
second continuation within each such block. The flat-source augmentation
square uses a single target augmentation, not one per outer coefficient;
its inner maps act on complete continuation sums. In the reverse-function
instance, actual native homs provide all four arrows to the binary theorem.
The hom round trips to the existing total-token presentation preserve global
augmentation and composition; they are not a compatibility theorem for
every contextad operation.
\paragraph{Nullary discriminator.}
Natural invertible absorption on both sides and natural invertible right
unit imply that an arrow $0\to I$ exists precisely when $0$ is initial.
The proof does not stipulate initiality. A parity-labelled fixture in the
other regime has two native maps with identical noninjective carrier,
total map and global augmentation. Their distinct maps on an unhit
coefficient survive an additional continuation step. The empty occurrence
list separately has multiplicative unit coefficient; a missing inner or
outer branch contributes the additive unit instead.
\paragraph{Boundary of transport.}
This is not the full native Fubini theorem under every Laplaza rig.
The missing work includes: (1) deriving the local diagrams from a chosen
standard rig presentation; (2) relating the flat native fibre and both
original unnormalized evaluation routes to the common ordered branch
presentation with all intervening transports and partitions; (3) integrating
the component into the dependent functor on arbitrary pairs of native arrows;
(4) the all-arity augmented comparison and its coherences for dependent
associativity, independent regrouping, permutations and units. The
arbitrary-length natural shuffle and its inverse concern the coefficient
tensor, not all four claims. These obligations are not extra assumptions of
the binary theorem. No full-project verification is asserted.
\end{document}
